\chapter{Introduction}
BLDC and PMSM motors are increasingly more popular in varous applications, such as robotics, drones, electric vehicles and many more, due to their high efficiency, 
ability to produce large torque at low speeds and high position accuracy. Unlike brushed DC or universal motors, 
BLDC motors need specialized driver and control to achieve commutation.

While commercial BLDC drivers are available on todays market, such as ODrive Micro, moteus-c1, 
they are often expensive in range of 70-90 USD per driver. \cite{moteus_c1} \cite{odrive_micro}
Therefore, the Department of Control Engineering at the Czech Technical University in Prague has developed an 
integrated three-phase motor driver called uFOC (micro FOC) with an STM-32 type microcontroller, 
which is a low-cost alternative to commercial drivers at BOM price of approximately 20 USD. \cite{uFOC}

The goal of the uFOC project was to develop a compact integrated motor driver for dynamic servo actuators, 
containing all components required for closed-loop torque control on a single printed circuit board.
The concept of tightly integrating the motor, driver electronics, 
and communication interfaces into a compact actuator unit was inspired by the modular actuator design presented
by Benjamin Katz in his work on low-cost actuators for dynamic robots. \cite{ben_katz}
The current revision of the uFOC driver has dimensions of 30×30~mm, approaching the practical physical limits imposed by components such as connectors.
The driver is intended for direct mounting onto the motor assembly, minimizing external wiring requirements to only power and CAN bus connections.
Two CAN connectors are provided to support daisy-chain communication in multi-actuator systems. \cite{uFOC}

To achive an efficient commutation and generate torque, the driver needs to implement a closed loop control.
The most common control method for BLDC motors is Field Oriented Control (FOC), 
which allows for precise control of the motor's torque, speed and position by controlling the current in the stator windings.

Goal of this thesis is to develop embedded firmware that implements FOC current control, speed and position control loops and communication via CAN bus for the uFOC driver. 
Additionally a Python client side example for CAN communication with the driver is provided.


